\documentclass[a4paper,11pt]{article}

\usepackage{amsmath,amssymb,mathtools,bm}
\usepackage{geometry}
\usepackage{booktabs}
\usepackage{xcolor}
\usepackage{hyperref}
\usepackage{microtype}
\geometry{top=23mm,bottom=25mm,left=24mm,right=24mm}
\hypersetup{colorlinks=true,linkcolor=blue!45!black,citecolor=blue!45!black,urlcolor=blue!50!black}
\setlength{\parindent}{1.2em}
\setlength{\parskip}{0.08em}
\numberwithin{equation}{section}

\newcommand{\dd}{\mathrm d}
\newcommand{\tr}{\operatorname{tr}}
\newcommand{\Lag}{\mathcal L}
\newcommand{\G}{\mathcal G}
\newcommand{\Id}{\mathbf 1}
\newcommand{\eps}{\varepsilon}

\begin{document}

\hypersetup{pageanchor=false}
\begin{titlepage}
\thispagestyle{empty}
\begin{center}
\vspace*{22mm}
{\Large\bfseries Exact stress-tensor flow in four-dimensional\\[2mm]
single-scalar Dirac--Born--Infeld theory\par}
\vspace{16mm}
{\large Ryosuke Suda\footnotemark\par}
\vspace{4mm}
{\normalsize
Graduate School of Science and Engineering, Saitama University,\\
255 Shimo-Okubo, Sakura-ku, Saitama 338-8570, Japan\par}
\vspace{18mm}
\end{center}

\noindent\textbf{Abstract.}
We consider four-dimensional Dirac--Born--Infeld theory containing one Abelian gauge field and one real scalar field, and derive a flow equation that expresses the derivative of the Lagrangian with respect to the deformation parameter solely in terms of local invariants of the stress tensor. For a non-null scalar gradient, the flow admits a closed rational expression. The same expression follows both from a direct calculation in the four-dimensional DBI theory and from dimensional reduction of the known five-dimensional Born--Infeld flow. When the scalar gradient is null, the rational expression degenerates and is replaced by a finite stress-tensor formula; in the limit where the scalar is removed, this formula reduces to the known four-dimensional Born--Infeld flow. We also derive an algebraic relation among the stress-tensor invariants realized by the DBI theory. The closed form was first suggested by a neural-network-assisted numerical search and is established analytically in the paper.

\vspace{8mm}
\noindent\textbf{Keywords:} Dirac--Born--Infeld theory; non-linear electrodynamics; stress-tensor deformation; D-branes

\vfill
\footnotetext{E-mail: \texttt{r.suda.813 [at] ms.saitama-u.ac.jp}}
\end{titlepage}

\hypersetup{pageanchor=true}
\setcounter{page}{1}

\section{Introduction}

Non-linear electrodynamics provides a broad setting in which to study electromagnetic duality, causality, and wave propagation beyond Maxwell theory. Born--Infeld (BI) theory is the canonical example. More recently, the structure of four-dimensional non-linear electrodynamics (NLED) has received renewed attention through the ModMax theory, which preserves both conformal invariance and electromagnetic duality, its Born--Infeld-type extensions, and classifications of theories without birefringence \cite{BornInfeld1934,Bandos2020,Sorokin2021,Ferko2023Birefringence,Ferko2023Duality}.

In parallel, deformations generated by local composite operators built from the stress tensor have been studied extensively. Their prototype is the two-dimensional $T\overline T$ deformation, an irrelevant deformation that nevertheless retains substantial analytic control and, in suitable cases, preserves integrability and allows exact access to finite-volume spectra \cite{SmirnovZamolodchikov2016}. There is no unique higher-dimensional analogue of $T\overline T$, but flow equations in which the variation of a Lagrangian with respect to a deformation parameter is expressed in terms of stress-tensor invariants have been identified for BI theory, ModMax--BI theory, and brane actions, among other examples \cite{Taylor2018,Babaei2022,Ferko2022,Conti2022,Ferko2023Duality,Brizio2026}.

The situation is more involved for Dirac--Born--Infeld (DBI) theories containing both a gauge field and worldvolume scalars. In generic dimension, the derivative of the DBI Lagrangian with respect to the deformation parameter contains invariants built explicitly from the field strength and scalar gradients, and these quantities need not be reconstructible from the stress tensor alone. Blair, Lahnsteiner, Obers and Yan studied flows generating D$p$-brane DBI actions and obtained complete flow equations including the gauge field for $p=1,2$, while for the four-dimensional worldvolume case $p=3$ they gave the result only to first order in the deformation parameter \cite{Blair2025}. Brizio, Kade, Sfondrini and Sorokin derived stress-tensor flows for BI theory in general dimension and analyzed DBI theories in two and three dimensions in detail \cite{Brizio2026}. They identified the extension to $d\geq4$ DBI theories as an open problem and pointed out that the four-dimensional theory with a single real scalar is a particularly promising case in which a formulation entirely in terms of stress-tensor invariants may exist. Relatedly, Ferko, Hu, Huang, Koutrolikos and Tartaglino-Mazzucchelli constructed stress-tensor-like flows for three-dimensional Born--Infeld theory and for the scalar sector of DBI theory, and identified the construction of a flow for the full gauge-plus-scalar DBI action as a natural open direction \cite{Ferko2024DBI}.

In this paper we address this four-dimensional single-scalar DBI theory. Before carrying out the analytic calculation, we numerically regressed the derivative of the DBI Lagrangian with respect to the deformation parameter against four invariants constructed from the stress tensor on the same field configurations. A neural network reproduced this quantity on independent test data with relative $L^2$ error $1.37\times10^{-2}$, indicating that the derivative may be a function of the stress-tensor invariants alone. A simple rational ansatz applied to independent numerical data then produced a candidate closed form.

Guided by this numerical observation, we derive the flow analytically in four dimensions. Intermediate expressions contain invariants built from the field strength and the scalar gradient, but identities specific to four-dimensional antisymmetric matrices allow these quantities to be eliminated in favor of stress-tensor invariants. We then derive the same four-dimensional flow by reducing the known five-dimensional Born--Infeld flow. The formula obtained for a non-null scalar gradient becomes indeterminate when the scalar gradient is null; we therefore treat this locus separately and recover the pure Born--Infeld flow as a limiting case. The details of the numerical search are collected in appendix~\ref{app:numerical-search}; none of the analytic derivations below assumes the numerically inferred formula.

All flow equations considered below are classical Lagrangian flow equations; no quantum definition of the four-dimensional deformation is assumed.

The paper is organized as follows. Section~\ref{sec:dbi} fixes the conventions for the four-dimensional DBI theory and its stress tensor. In section~\ref{sec:flow} we derive the four-dimensional flow equation directly and compare its weak-deformation limit with the perturbative $p=3$ result of ref.~\cite{Blair2025}. Section~\ref{sec:reduction} derives the same relation by dimensional reduction from five-dimensional Born--Infeld theory. Section~\ref{sec:null} treats null scalar gradients and the Born--Infeld limit. In section~\ref{sec:constraint} we derive an algebraic relation among the stress-tensor invariants realized by the DBI theory. Section~\ref{sec:discussion} summarizes the results and discusses their interpretation and possible extensions.

\section{Four-dimensional DBI theory}
\label{sec:dbi}

We first fix conventions for four-dimensional DBI theory with a single real scalar. This section introduces the action and the Hilbert stress tensor. Since the characteristic polynomial of the stress tensor will be used below, we distinguish throughout between the covariant tensor $T_{\mu\nu}$ and the mixed tensor $T^\mu{}_{\nu}$.

\subsection{Conventions and action}

We use metric signature $(-,+,+,+)$. In canonical four-dimensional normalization the deformation parameter has mass dimension $[\lambda]=-4$, and we take $\lambda>0$ in order to describe the real DBI branch. The algebraic formulas below may be analytically continued wherever the relevant square roots remain defined. The Abelian gauge potential is denoted by $A_\mu$, with field strength
\begin{equation}
 F_{\mu\nu}=\partial_\mu A_\nu-\partial_\nu A_\mu.
 \label{eq:F-definition}
\end{equation}
We denote the real scalar field by $\phi$. The four-dimensional DBI action is
\begin{equation}
 I_\lambda=\int\dd^4x\sqrt{-g}\,\Lag_\lambda,
 \qquad
 \Lag_\lambda=\frac{1-\Delta}{\lambda},
 \label{eq:action}
\end{equation}
where
\begin{equation}
 \G_{\mu\nu}
 =g_{\mu\nu}+\lambda\,\partial_\mu\phi\,\partial_\nu\phi
 +\sqrt\lambda F_{\mu\nu},
 \qquad
 \Delta=\sqrt{\frac{\det\G}{\det g}}.
 \label{eq:dbi-matrix}
\end{equation}
The branch is chosen so that $\Delta\to1$ as $\lambda\to0$.

We write the Lorentz norm of the scalar gradient as
\begin{equation}
 X=\partial_\mu\phi\,\partial^\mu\phi.
 \label{eq:X-definition}
\end{equation}
For later use, we also introduce the standard NLED invariant
\begin{equation}
 S=-\frac14F_{\mu\nu}F^{\mu\nu}.
 \label{eq:S-definition}
\end{equation}
Expanding eq.~\eqref{eq:dbi-matrix} around $\lambda=0$ gives
\begin{equation}
 \Lag_\lambda=S-\frac12X+O(\lambda).
 \label{eq:seed}
\end{equation}
Thus the seed theory at $\lambda=0$ consists of a Maxwell field and a free massless scalar.

\subsection{Stress tensor}

To formulate the flow in terms of stress-tensor invariants, we first determine the Hilbert stress tensor of the DBI action. We define
\begin{equation}
 T_{\mu\nu}
 =-\frac{2}{\sqrt{-g}}\frac{\delta I_\lambda}{\delta g^{\mu\nu}}.
 \label{eq:stress-definition}
\end{equation}
Metric variation of eq.~\eqref{eq:action} gives
\begin{equation}
 T^{\mu\nu}
 =\frac1\lambda
 \left[g^{\mu\nu}-\Delta(\G^{-1})^{(\mu\nu)}\right],
 \qquad
 (\G^{-1})^{(\mu\nu)}
 =\frac12\left[(\G^{-1})^{\mu\nu}+(\G^{-1})^{\nu\mu}\right].
 \label{eq:stress-dbi}
\end{equation}
All matrix invariants of the stress tensor below are taken for the mixed tensor
\begin{equation}
 T^\mu{}_{\nu}=g^{\mu\rho}T_{\rho\nu}.
 \label{eq:mixed-T}
\end{equation}
Equation~\eqref{eq:stress-dbi} will be the starting point for the calculation below. In the next section we compute the derivative of the DBI Lagrangian on the same field configuration and eliminate the explicit field variables in favor of invariants of $T^\mu{}_{\nu}$.

\section{Four-dimensional flow equation}
\label{sec:flow}

In this section we differentiate the four-dimensional DBI Lagrangian with respect to the deformation parameter while holding the metric, field strength, and scalar gradient fixed, and then express the result in terms of stress-tensor invariants. We first introduce the relevant invariants and briefly describe the numerical search that motivated the analytic calculation. We then derive the flow directly from the DBI Lagrangian without assuming the numerically inferred expression.

We define
\begin{equation}
 Q_\lambda\equiv
 \left.\frac{\partial\Lag_\lambda}{\partial\lambda}\right|_{g_{\mu\nu},\,F_{\mu\nu},\,\partial_\mu\phi}.
 \label{eq:Q-definition}
\end{equation}
The goal is to express $Q_\lambda$ solely through local invariants of the stress tensor. Introduce
\begin{equation}
 M^\mu{}_{\nu}=\delta^\mu{}_{\nu}-\lambda T^\mu{}_{\nu}.
 \label{eq:M-definition}
\end{equation}
For a four-dimensional matrix $N$, define the elementary symmetric polynomials $e_k(N)$ by
\begin{equation}
 \det(z\Id-N)
 =z^4-e_1(N)z^3+e_2(N)z^2-e_3(N)z+e_4(N),
 \label{eq:elementary-definition}
\end{equation}
and write $E_k=e_k(M)$. In particular, for the stress tensor itself,
\begin{align}
 e_2(T)&=\frac12\left[(\tr T)^2-\tr(T^2)\right],\nonumber\\
 e_3(T)&=\frac16\left[(\tr T)^3-3\tr T\,\tr(T^2)+2\tr(T^3)\right].
 \label{eq:e2e3-explicit}
\end{align}

The numerical search was carried out in the region where $X\neq0$ and the scalar gradient is timelike. Random DBI field configurations were used to compute $(E_1,E_2,E_3,E_4)$ and $\lambda^2Q_\lambda$, and a neural network was trained to regress the map
\begin{equation}
 (E_1,E_2,E_3,E_4)\longmapsto \lambda^2Q_\lambda.
 \label{eq:numerical-clue}
\end{equation}
Its small error on independent test data suggested that $Q_\lambda$ could be written as a function of these invariants alone. The network architecture, loss function, and rational readout are described in appendix~\ref{app:numerical-search}. We now proceed analytically and do not assume this numerical result.

For an antisymmetric $4\times4$ matrix $F^\mu{}_{\nu}$, the Cayley--Hamilton identity reduces to a form involving the two standard four-dimensional electromagnetic invariants. In addition to $S$, define
\begin{equation}
 P=-\frac14F_{\mu\nu}\widetilde F^{\mu\nu},
 \qquad
 \widetilde F^{\mu\nu}
 =\frac{1}{2\sqrt{-g}}\eps^{\mu\nu\rho\sigma}F_{\rho\sigma},
 \qquad \eps^{0123}=+1.
 \label{eq:P-definition}
\end{equation}
Then
\begin{equation}
 (\sqrt\lambda F)^4-2\lambda S(\sqrt\lambda F)^2-\lambda^2P^2\Id=0.
 \label{eq:F-characteristic-main}
\end{equation}
We will also use the matrix determinant lemma
\begin{equation}
 \det(A+uv^{\mathsf T})=\det A\left(1+v^{\mathsf T}A^{-1}u\right)
 \label{eq:rank-one-det-main}
\end{equation}
for an invertible matrix $A$ and column vectors $u,v$.

Applying these relations to eq.~\eqref{eq:stress-dbi} gives first
\begin{equation}
 E_4=\det M=1+\lambda X.
 \label{eq:E4}
\end{equation}
The intermediate expressions for $E_1$ and $E_3$ involve the mixed gauge--scalar invariant
\begin{equation}
 Z=\partial_\mu\phi\,F^\mu{}_{\rho}F^\rho{}_{\nu}\,\partial^\nu\phi,
 \label{eq:Z-definition}
\end{equation}
which does not appear in the final flow equation. A direct application of eqs.~\eqref{eq:F-characteristic-main} and \eqref{eq:rank-one-det-main} to the DBI matrix yields
\begin{align}
 \Delta^2
 &=1+\lambda X-2\lambda S
 +\lambda^2\left(Z-2SX-P^2\right),
 \label{eq:delta-fields}\\
 \Delta E_1
 &=4+3\lambda X-4\lambda S+\lambda^2(Z-2SX),
 \label{eq:E1-fields}\\
 \Delta E_3
 &=(1+\lambda X)(4+\lambda X-4\lambda S)+\lambda^2Z.
 \label{eq:E3-fields}
\end{align}
The matrix algebra leading to these formulas is given in appendix~\ref{app:matrix-identities}.

On the other hand, eq.~\eqref{eq:action} implies
\begin{equation}
 \lambda^2Q_\lambda
 =\Delta-1-\lambda\frac{\partial\Delta}{\partial\lambda}.
 \label{eq:Q-delta}
\end{equation}
Differentiating eq.~\eqref{eq:delta-fields} gives
\begin{equation}
 2\Delta\left(1+\lambda^2Q_\lambda\right)
 =2+\lambda X-2\lambda S.
 \label{eq:Q-field-relation}
\end{equation}
Taking the difference of eqs.~\eqref{eq:E1-fields} and \eqref{eq:E3-fields} and using eq.~\eqref{eq:E4}, we obtain
\begin{equation}
 E_1-E_3
 =2(1-E_4)\left(1+\lambda^2Q_\lambda\right).
 \label{eq:undivided-flow}
\end{equation}
This relation remains valid at $X=0$ and will also be used in section~\ref{sec:null}.

For $X\neq0$, eq.~\eqref{eq:E4} implies $E_4\neq1$, so eq.~\eqref{eq:undivided-flow} can be solved for $Q_\lambda$:
\begin{equation}
 \lambda^2Q_\lambda
 =\frac{E_1-E_3+2E_4-2}{2(1-E_4)}.
 \label{eq:flow-E}
\end{equation}
To rewrite this entirely in terms of $T^\mu{}_{\nu}$, use $M=\Id-\lambda T$. With $e_1(T)=\tr T$ and $e_4(T)=\det T$,
\begin{align}
 E_1&=4-\lambda\tr T,\nonumber\\
 E_3&=4-3\lambda\tr T+2\lambda^2e_2(T)-\lambda^3e_3(T),\nonumber\\
 E_4&=1-\lambda\tr T+\lambda^2e_2(T)-\lambda^3e_3(T)+\lambda^4\det T.
 \label{eq:E-to-T}
\end{align}
It follows that
\begin{equation}
 Q_\lambda
 =\frac{2\lambda\det T-e_3(T)}
 {2\left[\tr T-\lambda e_2(T)+\lambda^2e_3(T)-\lambda^3\det T\right]}.
 \label{eq:main-flow}
\end{equation}
Moreover, comparing eqs.~\eqref{eq:E4} and \eqref{eq:E-to-T} gives
\begin{equation}
 \tr T-\lambda e_2(T)+\lambda^2e_3(T)-\lambda^3\det T=-X.
 \label{eq:denominator-X}
\end{equation}
Thus eq.~\eqref{eq:main-flow} applies for $X\neq0$, while the locus $X=0$ must be treated separately.

Equation~\eqref{eq:main-flow} should be understood as an exact differential identity along the one-parameter DBI family. It gives one stress-tensor-only representative of the deformation generator on the locus of stress tensors realized by DBI configurations. As shown in section~\ref{sec:constraint}, that locus obeys an algebraic constraint, so extensions of the same operator away from the DBI locus are not unique. We do not prove here that eq.~\eqref{eq:main-flow}, regarded as an autonomous equation on a larger space of Lagrangians with the Maxwell--scalar seed, has a unique solution equal to DBI.

As a first consistency check, consider the weak-deformation limit. For a seed configuration with $X\neq0$, eq.~\eqref{eq:main-flow} gives
\begin{equation}
 Q_0=-\frac{e_3(T_0)}{2\tr T_0},
 \label{eq:leading-flow}
\end{equation}
where $T_0$ is the stress tensor of the Maxwell--scalar theory in eq.~\eqref{eq:seed}.

Blair, Lahnsteiner, Obers and Yan denote their deformation parameter by $t$ \cite{Blair2025}. Restricting their flat-worldvolume DBI Lagrangian to $p=3$ and one real scalar, and identifying $t=\lambda$ and $X^{A'}=\phi$, gives the normalization in eq.~\eqref{eq:action} after extracting the overall factor $t^{-1}$ from the four-dimensional determinant. In these conventions, the $t^0$ term in their $p=3$ flow can be written as
\begin{equation}
 Q_0
 =\frac14\tr(T_0^2)-\frac18(\tr T_0)^2
 -\frac18\left[\tr(F^4)-\frac14\bigl(\tr F^2\bigr)^2\right],
 \label{eq:blair-leading}
\end{equation}
where $F$ denotes the mixed matrix $F^\mu{}_{\nu}$. The comparison can be made explicitly. For the Maxwell--scalar seed one finds
\begin{align}
 \tr T_0&=-X,\nonumber\\
 \tr(T_0^2)&=4(S^2+P^2)+X^2+2SX-2Z,\nonumber\\
 e_3(T_0)&=X\left(S^2+P^2+SX-Z+\frac{X^2}{4}\right),\nonumber\\
 \tr(F^4)-\frac14\bigl(\tr F^2\bigr)^2&=4(S^2+P^2).
 \label{eq:seed-identities}
\end{align}
For $X\neq0$, both eqs.~\eqref{eq:leading-flow} and \eqref{eq:blair-leading} therefore give
\begin{equation}
 Q_0=\frac12\left(S^2+P^2+SX-Z+\frac{X^2}{4}\right).
 \label{eq:leading-flow-fields}
\end{equation}
At $X=0$ the ratio in eq.~\eqref{eq:leading-flow} is indeterminate, but eq.~\eqref{eq:blair-leading} remains finite and will agree with the null-branch expansion in section~\ref{sec:null}.

There is also a complementary one-field truncation. Setting $F_{\mu\nu}=0$ gives the scalar DBI (Dirac) theory
\begin{equation}
 \Lag_\lambda^{(\phi)}=\frac{1-\Delta_\phi}{\lambda},
 \qquad
 \Delta_\phi=\sqrt{1+\lambda X},
 \label{eq:pure-scalar-L}
\end{equation}
with mixed stress tensor
\begin{equation}
 T^\mu{}_{\nu}=\Lag_\lambda^{(\phi)}\delta^\mu{}_{\nu}
 +\frac{1}{\Delta_\phi}\,\partial^\mu\phi\,\partial_\nu\phi.
 \label{eq:pure-scalar-T}
\end{equation}
Substitution into eq.~\eqref{eq:main-flow} yields
\begin{equation}
 Q_\lambda^{(\phi)}
 =\frac{(\Delta_\phi-1)^2}{2\lambda^2\Delta_\phi}.
 \label{eq:pure-scalar-Q}
\end{equation}
Ferko et al.~\cite{Ferko2024DBI} use $x_1=X/2$ and show that their stress-tensor operator for the scalar sector reduces to
\begin{equation}
 \frac{\partial\Lag}{\partial\lambda}
 =-\frac12\Lag^2+\Lag x_1\frac{\partial\Lag}{\partial x_1}.
 \label{eq:ferko-scalar-flow}
\end{equation}
For eq.~\eqref{eq:pure-scalar-L}, $\partial\Lag/\partial x_1=-1/\Delta_\phi$, and the right-hand side of eq.~\eqref{eq:ferko-scalar-flow} is exactly eq.~\eqref{eq:pure-scalar-Q}. Thus the present gauge-plus-scalar result reproduces the known scalar DBI flow when the gauge field is removed. The opposite truncation, $\partial_\mu\phi=0$, is recovered from the $X=0$ formula in section~\ref{sec:null} and gives the known four-dimensional Born--Infeld trace flow.

Equation~\eqref{eq:main-flow} therefore resums the deformation parameter to all orders, while reproducing the known perturbative $p=3$ result at weak deformation and both known one-field truncations. The derivation above uses only the four-dimensional DBI Lagrangian. Since the DBI action itself arises by dimensional reduction of higher-dimensional Born--Infeld theory, it is natural to ask whether the flow equation behaves in the same way. We address this next.

\section{Five-dimensional Born--Infeld reduction}
\label{sec:reduction}

The previous section derived the flow directly in four dimensions. We now obtain the same expression from the known five-dimensional Born--Infeld flow. A gauge-potential component along a compact direction becomes a scalar field after reduction, so reducing five-dimensional Born--Infeld theory on one direction produces precisely the four-dimensional DBI system considered above. We first verify the reduction at the level of the action and then apply it to the flow equation.

Five-dimensional quantities will carry hats. Let the compact coordinate be $y\sim y+L_y$. We do not consider a general Kaluza--Klein reduction: throughout this section we take a product background with $\widehat g_{\mu y}=0$ and $\widehat g_{yy}=1$, omit fluctuations of the Kaluza--Klein vector and radion, and retain only zero modes, so all fields are independent of $y$.

It is useful to distinguish the canonically normalized deformation parameters before and after reduction. Denote the five-dimensional parameter by $\widehat\lambda$ and define the four-dimensional zero modes by
\begin{equation}
 \lambda=\frac{\widehat\lambda}{L_y},
 \qquad
 A_\mu=\sqrt{L_y}\,\widehat A_\mu,
 \qquad
 \phi=\sqrt{L_y}\,\widehat A_y.
 \label{eq:reduction-rescaling}
\end{equation}
Thus $[\widehat\lambda]=-5$ and $[\lambda]=-4$ in canonical normalization. The field strengths obey
$\widehat F_{\mu\nu}=F_{\mu\nu}/\sqrt{L_y}$ and
$\widehat F_{\mu y}=\partial_\mu\phi/\sqrt{L_y}$.

We take
\begin{equation}
 \widehat g_{MN}
 =\begin{pmatrix}g_{\mu\nu}&0\\0&1\end{pmatrix},
 \qquad
 \widehat\G_{MN}=\widehat g_{MN}+\sqrt{\widehat\lambda}\,\widehat F_{MN}.
 \label{eq:5d-G}
\end{equation}
The five-dimensional action and determinant factor are
\begin{equation}
 \widehat I_{\widehat\lambda}=\int\dd^5x\sqrt{-\widehat g}\,
 \widehat\Lag_{\widehat\lambda},
 \qquad
 \widehat\Lag_{\widehat\lambda}=\frac{1-\widehat\Delta}{\widehat\lambda},
 \qquad
 \widehat\Delta=\sqrt{\frac{\det\widehat\G}{\det\widehat g}},
 \label{eq:5d-action}
\end{equation}
and the Hilbert stress tensor is
\begin{equation}
 \widehat T_{MN}=-\frac{2}{\sqrt{-\widehat g}}
 \frac{\delta\widehat I_{\widehat\lambda}}{\delta\widehat g^{MN}},
 \qquad
 \widehat T^M{}_{N}=\widehat g^{MP}\widehat T_{PN}.
 \label{eq:5d-stress}
\end{equation}
Using eq.~\eqref{eq:reduction-rescaling}, the zero-mode matrix takes the form
\begin{equation}
 \widehat\G
 =\begin{pmatrix}
 g+\sqrt\lambda F&\sqrt\lambda\,\partial\phi\\
 -\sqrt\lambda\,(\partial\phi)^{\mathsf T}&1
 \end{pmatrix}.
 \label{eq:5d-block}
\end{equation}
Since the lower-right block is unity, we may use
\begin{equation}
 \det\begin{pmatrix}A&b\\ c^{\mathsf T}&1\end{pmatrix}
 =\det\left(A-bc^{\mathsf T}\right).
 \label{eq:schur-explicit}
\end{equation}
Applied to eq.~\eqref{eq:5d-block}, this gives
\begin{equation}
 \det\widehat\G
 =\det\left(g+\sqrt\lambda F
 +\lambda\,\partial\phi\,(\partial\phi)^{\mathsf T}\right)
 =\det\G.
 \label{eq:schur-action}
\end{equation}
After integrating over $y$, the factor $L_y$ combines with
$\widehat\lambda=L_y\lambda$, so that
$\widehat I_{\widehat\lambda}=I_\lambda$ with the four-dimensional action
in eq.~\eqref{eq:action}. Metric variation similarly gives
$T^\mu{}_{\nu}=L_y\,\widehat T^\mu{}_{\nu}$ for the four-dimensional block, and hence
$\widehat\lambda\,\widehat T^\mu{}_{\nu}=\lambda T^\mu{}_{\nu}$.

The inverse matrix needed for reducing the flow follows from the same Schur complement:
\begin{equation}
 \widehat\G^{-1}
 =\begin{pmatrix}
 \G^{-1} & -\sqrt\lambda\,\G^{-1}\partial\phi\\
 \sqrt\lambda\,(\partial\phi)^{\mathsf T}\G^{-1}
 &1-\lambda(\partial\phi)^{\mathsf T}\G^{-1}\partial\phi
 \end{pmatrix}.
 \label{eq:5d-block-inverse}
\end{equation}
As in four dimensions, define
\begin{equation}
 \widehat M^M{}_{N}\equiv\delta^M{}_{N}
 -\widehat\lambda\,\widehat T^M{}_{N}
 =\widehat\Delta\,(\widehat\G^{-1})^{(MP)}\widehat g_{PN}.
 \label{eq:Mhat-definition}
\end{equation}
Equation~\eqref{eq:schur-action} implies $\widehat\Delta=\Delta$, so the trace and determinant of $\widehat M$ can be expressed directly through four-dimensional quantities.

The known five-dimensional Born--Infeld flow is \cite{Brizio2026}
\begin{equation}
 \frac{\partial\widehat\Lag_{\widehat\lambda}}{\partial\widehat\lambda}
 =\frac{1}{4\widehat\lambda^2}
 -\frac{1}{4\widehat\lambda}\tr\widehat T
 -\frac{1}{4\widehat\lambda^2}
 \det(\Id_5-\widehat\lambda\widehat T).
 \label{eq:5d-flow}
\end{equation}
Since $\Lag_\lambda=L_y\widehat\Lag_{\widehat\lambda}$ and
$\widehat\lambda=L_y\lambda$ at fixed $L_y$ and fixed four-dimensional fields,
\begin{equation}
 \lambda^2 Q_\lambda
 =\widehat\lambda^2
 \frac{\partial\widehat\Lag_{\widehat\lambda}}{\partial\widehat\lambda}.
 \label{eq:dimensionless-flow-reduction}
\end{equation}
Using eqs.~\eqref{eq:5d-block-inverse}, \eqref{eq:Mhat-definition}, and the four-dimensional stress tensor in eq.~\eqref{eq:stress-dbi}, one finds
\begin{align}
 \det(\Id_5-\widehat\lambda\widehat T)&=\Delta,
 \label{eq:det-Mhat}\\
 \tr(\Id_5-\widehat\lambda\widehat T)
 &=E_1+\Delta\left[1-
 \lambda\,\partial_\mu\phi(\G^{-1})^{(\mu\nu)}\partial_\nu\phi\right].
 \label{eq:trace-Mhat}
\end{align}
Equations~\eqref{eq:5d-flow}--\eqref{eq:trace-Mhat} then reduce to
\begin{equation}
 \lambda^2Q_\lambda
 =\frac14\left[E_1-4-
 \Delta\lambda\,\partial_\mu\phi(\G^{-1})^{(\mu\nu)}\partial_\nu\phi\right].
 \label{eq:reduced-before-elimination}
\end{equation}
The remaining scalar contraction can be eliminated in favor of four-dimensional stress-tensor invariants. Combining eqs.~\eqref{eq:E1-fields}, \eqref{eq:E3-fields}, and \eqref{eq:E4} gives
\begin{equation}
 (E_4-1)\Delta\lambda\,
 \partial_\mu\phi(\G^{-1})^{(\mu\nu)}\partial_\nu\phi
 =E_1(1+E_4)-2E_3.
 \label{eq:elimination}
\end{equation}
For $X\neq0$, substituting eq.~\eqref{eq:elimination} into eq.~\eqref{eq:reduced-before-elimination} reproduces eq.~\eqref{eq:flow-E}. This derivation does not differentiate the four-dimensional DBI Lagrangian with respect to $\lambda$ and therefore provides a second route to the same four-dimensional flow from the higher-dimensional Born--Infeld relation. The final elimination step does reuse the four-dimensional algebraic identities for $E_1$, $E_3$, and $E_4$, so this route is a complementary derivation and consistency check rather than a completely independent proof of all the matrix algebra.

Both derivations divide by $1-E_4$ in the final step. According to eq.~\eqref{eq:E4}, this factor vanishes precisely at $X=0$. We therefore analyze this locus separately.

\section{Null scalar gradient}
\label{sec:null}

The expression in eq.~\eqref{eq:flow-E} was derived under the assumption $X\neq0$. Equation~\eqref{eq:E4} shows that $X=0$ is equivalent to $E_4=1$, in which case both numerator and denominator of eq.~\eqref{eq:flow-E} vanish. The locus $X=0$ contains both a nonzero null scalar gradient and the special case $\partial_\mu\phi=0$. To obtain a finite expression there, we return to the undivided relation in eq.~\eqref{eq:undivided-flow} and use the characteristic polynomial of $M$.

For four-dimensional DBI theory with a single real scalar, the characteristic polynomial of $M$ factorizes into two quadratics, as shown in appendix~\ref{app:matrix-identities}. On the locus $X=0$, this factorization gives
\begin{equation}
 \lambda^2Q_\lambda
 =\frac14\left[E_1+\sqrt{E_1^2-4E_2+8}\right]-1.
 \label{eq:null-E}
\end{equation}
The sign can be fixed algebraically, without referring only to a path of approach from $X\neq0$. On $X=0$, eqs.~\eqref{eq:E1-fields}, \eqref{eq:Q-field-relation}, and the factorization coefficients of appendix~\ref{app:matrix-identities} give
\begin{equation}
 \beta=\frac{2(1-\lambda S)}{\Delta},
 \qquad
 \alpha=\frac{2-2\lambda S+\lambda^2Z}{\Delta},
 \qquad
 \beta-\alpha=-\frac{\lambda^2Z}{\Delta}.
 \label{eq:null-alpha-beta-fields}
\end{equation}
For a nonzero null gradient $v_\mu=\partial_\mu\phi$, define $w_\mu=F_{\mu\nu}v^\nu$. Antisymmetry gives $w\cdot v=0$, and a vector orthogonal to a nonzero null vector is spacelike or null. Hence $w^2\geq0$ and
\begin{equation}
 Z=-w_\mu w^\mu\leq0.
 \label{eq:null-Z-sign}
\end{equation}
On the positive-$\Delta$ DBI branch, eq.~\eqref{eq:null-alpha-beta-fields} therefore implies $\beta\geq\alpha$, selecting the plus sign in eq.~\eqref{eq:null-E}. If $v_\mu=0$, then $Z=0$ and the two roots coincide.

Expanding the null formula at weak deformation gives
\begin{equation}
 Q_0\big|_{X=0}=\frac12\left(S^2+P^2-Z\right).
 \label{eq:null-weak-limit}
\end{equation}
This is precisely eq.~\eqref{eq:leading-flow-fields} at $X=0$ and also follows directly from the perturbative expression in eq.~\eqref{eq:blair-leading} using eq.~\eqref{eq:seed-identities}.

Rewriting $E_1$ and $E_2$ in terms of stress-tensor traces gives
\begin{equation}
 Q_\lambda
 =\frac{-\lambda\tr T+
 \sqrt{4\lambda\tr T+\lambda^2\left[2\tr(T^2)-(\tr T)^2\right]}}
 {4\lambda^2}.
 \label{eq:null-T}
\end{equation}

If the scalar gradient is further set to zero, the theory reduces to four-dimensional Born--Infeld electrodynamics. Equation~\eqref{eq:null-E} then simplifies to
\begin{equation}
 Q_\lambda=-\frac{1}{4\lambda}\tr T,
 \label{eq:pure-BI}
\end{equation}
in agreement with the known four-dimensional Born--Infeld flow \cite{Ferko2022,Brizio2026}. Thus a stress-tensor-only flow also exists at $X=0$, and the pure Born--Infeld result is recovered continuously.

Having treated both $X\neq0$ and $X=0$, we now examine the algebraic relations among the invariants $E_1,E_2,E_3,E_4$ realized by DBI configurations.

\section{Algebraic relation among stress-tensor invariants}
\label{sec:constraint}

The flow has been written in terms of $E_k=e_k(\Id-\lambda T)$. It is useful to determine whether these four quantities are independent when they arise from DBI field configurations. The factorization of the characteristic polynomial derived in appendix~\ref{app:matrix-identities} implies
\begin{align}
 (E_4-1)^2E_2={}&-E_1^2E_4+E_1E_3(E_4+1)-E_3^2\nonumber\\
 &+E_4^3-E_4^2-E_4+1.
 \label{eq:algebraic-relation}
\end{align}
Hence the four $E_k$ generated by DBI field configurations are not independent.

Equation~\eqref{eq:algebraic-relation} shows directly that $(E_1,E_2,E_3,E_4)$ cannot vary independently on DBI field configurations. The existence of the flow in eq.~\eqref{eq:main-flow} therefore does not follow from a simple counting argument based on the number of available invariants. The same relation also makes the algebraic representation of the flow operator non-unique: any multiple of the difference between the two sides of eq.~\eqref{eq:algebraic-relation} may be added to a stress-tensor expression without changing its value on DBI configurations.

\section{Summary and discussion}
\label{sec:discussion}

We have derived a stress-tensor flow for four-dimensional DBI theory containing one real scalar and one Abelian gauge field. For $X\neq0$, the derivative $\partial\Lag_\lambda/\partial\lambda$ is given by eq.~\eqref{eq:main-flow}, which depends only on local invariants of the mixed stress tensor. The same expression follows from a direct four-dimensional calculation and from dimensional reduction of the five-dimensional Born--Infeld flow. At $X=0$, eq.~\eqref{eq:null-T} provides a finite stress-tensor expression, and the pure Born--Infeld limit reduces to the known trace flow.

The immediate significance of the result is structural. Stress-tensor flows provide a way of describing trajectories in the space of classical Lagrangians without referring explicitly to all microscopic field invariants. In four-dimensional NLED this viewpoint has been useful for organizing families of duality-invariant theories and their deformations \cite{Ferko2023Duality,Ferko2023Birefringence}. The present calculation shows that such a description can persist when a gauge field and a worldvolume scalar are present simultaneously, at least in the special four-dimensional single-scalar DBI theory. In this sense it provides a four-dimensional single-scalar realization of the gauge-plus-scalar problem highlighted in ref.~\cite{Ferko2024DBI}. Since the undeformed limit is Maxwell theory plus a free massless scalar, one natural question is whether eq.~\eqref{eq:main-flow}, together with the seed Lagrangian in eq.~\eqref{eq:seed} and the $X=0$ prescription, uniquely reconstructs the DBI family within an appropriate class of local derivative-free Lagrangians. Establishing such a uniqueness statement would give an alternative characterization of this DBI theory directly in terms of its stress tensor.

The dimensional-reduction derivation provides a second direction. Brizio \textit{et al.} found stress-tensor flows for Born--Infeld theory in general dimension and for DBI theories in two and three dimensions \cite{Brizio2026}. The present five-to-four-dimensional reduction gives another explicit case in which a higher-dimensional Born--Infeld flow descends to a DBI flow, but only after the explicit field-strength dependence can be eliminated. This suggests studying dimensional reduction as a way to classify when a DBI deformation can be written solely in terms of stress-tensor invariants. The algebraic constraint in eq.~\eqref{eq:algebraic-relation} is relevant to such a classification: stress-tensor expressions that differ by a function vanishing on DBI configurations represent the same flow on this theory, so these algebraic relations must be taken into account when comparing candidate flow operators.

Several extensions are then concrete. With multiple worldvolume scalars, the rank-one structure used here is lost, while in other spacetime dimensions the algebra of a two-form is different; both changes can obstruct the elimination of explicit $F_{\mu\nu}$ invariants. It would therefore be useful to determine systematically which dimensions and field contents admit a stress-tensor-only description. Another question is whether the bosonic relation found here admits a supersymmetric completion or a formulation compatible with electromagnetic duality in the gauge sector, paralleling developments for four-dimensional NLED stress-tensor flows \cite{Ferko2023Birefringence,Ferko2023Duality}. Finally, the numerical search used in this work can serve as a diagnostic in these more complicated theories: it can test whether a candidate deformation is approximately single-valued as a function of stress-tensor invariants before one attempts an analytic elimination. The four-dimensional single-scalar theory provides a controlled test case in which this numerical strategy is known to lead to an exact relation.

All statements in this paper concern classical Lagrangian flows. Whether the relation admits a useful quantum interpretation in four dimensions is a separate question.

\section*{Acknowledgments}

\paragraph{Use of generative AI.}
OpenAI ChatGPT was used in this project to assist with exploration of research directions, preparation of numerical and symbolic-computation code, literature searches, and organization and editing of the manuscript. The analytic and numerical calculations and the references supporting the scientific claims were checked by the author, who takes responsibility for the content of the paper.

\section*{Data availability}

No external experimental or observational data were used in this work. Numerical validation data generated for the study are maintained in the accompanying repository.

\section*{Code availability}

Code for the numerical search, checks of the analytic formulas, and the five-to-four-dimensional reduction, together with numerical outputs, is maintained in the accompanying repository:\par\smallskip\noindent\url{https://github.com/ryos-physrockme/dbi-stress-tensor-flow}.\par\smallskip\noindent Code under the repository's \texttt{code/} directory is distributed under the MIT License.

\appendix

\section{Neural-network-assisted numerical search}
\label{app:numerical-search}

Before the analytic derivation, we investigated numerically whether $Q_\lambda$ could be predicted from stress-tensor invariants alone. On the flat metric $\eta=\mathrm{diag}(-1,1,1,1)$, define dimensionless fields
\begin{equation}
 a_\mu=\sqrt\lambda\,\partial_\mu\phi,
 \qquad
 b_{\mu\nu}=\sqrt\lambda\,F_{\mu\nu}.
 \label{eq:dimensionless-fields-nn}
\end{equation}
For a timelike scalar gradient, a Lorentz transformation brings $a_\mu$ to $(a,0,0,0)$, and the remaining spatial rotations reduce the independent nonzero components of $b_{\mu\nu}$ to three. We used
\begin{equation}
 b_{01}=e,\qquad b_{23}=p,\qquad b_{31}=b
\end{equation}
and drew samples uniformly from
\begin{equation}
 a\in[0.18,0.72],\quad
 e\in[-0.42,0.42],\quad
 p,b\in[-0.52,0.52].
 \label{eq:nn-sampling-range}
\end{equation}
We retained only configurations on the real DBI branch and excluded points with $\Delta<0.30$, $|1-E_4|<5\times10^{-3}$, or $|\lambda^2Q_\lambda|>2$ in order to stabilize the initial regression.

For each field configuration, the DBI Lagrangian and stress tensor were evaluated directly and used to form
\begin{equation}
 \bm{x}=(E_1,E_2,E_3,E_4)\in\mathbb R^4,
 \qquad
 y=\lambda^2Q_\lambda\in\mathbb R.
 \label{eq:nn-data}
\end{equation}
The neural network approximated the map
\begin{equation}
 f_\theta:\mathbb R^4\longrightarrow\mathbb R,
 \qquad
 (E_1,E_2,E_3,E_4)\longmapsto f_\theta(E_1,E_2,E_3,E_4)\simeq\lambda^2Q_\lambda.
 \label{eq:nn-map}
\end{equation}
Inputs and targets were standardized using the mean and standard deviation of the training set. If $\widetilde{\bm{x}}_n$ and $\widetilde y_n$ denote the standardized input and target of sample $n$, the loss was the mean-squared error
\begin{equation}
 \mathcal J(\theta)
 =\frac1{N_{\mathrm{train}}}
 \sum_{n\in\mathcal D_{\mathrm{train}}}
 \left[f_\theta(\widetilde{\bm{x}}_n)-\widetilde y_n\right]^2.
 \label{eq:nn-loss}
\end{equation}
No known flow equation or analytic candidate was included in the loss.

The network architecture was $4$--$64$--$64$--$64$--$1$, with $\tanh$ activations in the hidden layers. We generated $18{,}000$ field configurations, using $80\%$ for training and $20\%$ for testing. The network was trained with Adam, learning rate $2\times10^{-3}$, batch size $512$, and $260$ epochs. On the test set, the root-mean-square error was $7.31\times10^{-4}$ and the relative $L^2$ error was $1.37\times10^{-2}$. This regression accuracy by itself does not establish global single-valuedness, especially because the $E_k$ realized by DBI configurations lie on the algebraic hypersurface derived in section~\ref{sec:constraint}; its role was only to motivate the subsequent exact readout and analytic calculation.

To extract an analytic candidate, we generated an independent set of $12{,}000$ canonical-frame configurations and considered the simplest rational ansatz whose numerator and denominator are affine functions of the $E_k$,
\begin{equation}
 y=\frac{a_0+a_1E_1+a_2E_2+a_3E_3+a_4E_4}
 {b_0+b_1E_1+b_2E_2+b_3E_3+b_4E_4}.
 \label{eq:rational-ansatz}
\end{equation}
After clearing the denominator, the coefficients $a_i,b_i$ obey a homogeneous linear system. The two smallest singular values of the resulting design matrix were
\begin{equation}
 \sigma_{\min}=3.71\times10^{-14},
 \qquad
 \sigma_{\mathrm{next}}=2.41\times10^{-3},
 \label{eq:svd-gap}
\end{equation}
showing a clear one-dimensional numerical null direction. The associated singular vector reproduced, up to an overall normalization, the coefficients in eq.~\eqref{eq:flow-E}; the RMS residual of the affine-over-affine readout was $1.73\times10^{-14}$.

After obtaining this candidate, we tested eq.~\eqref{eq:flow-E} on $30{,}000$ field configurations generated without using the canonical frame, with all four components of $a_\mu$ and all six independent components of $b_{\mu\nu}$ sampled independently. The final claims of the paper are based on the analytic derivation in the main text and appendix~\ref{app:matrix-identities}, rather than on these numerical tests. The code and numerical outputs are available in the repository listed in the code-availability statement.

\section{Auxiliary matrix identities}
\label{app:matrix-identities}

\subsection{Determinant identities in four dimensions}

For compactness, define within this appendix
\begin{equation}
 C_\pm{}^\mu{}_{\nu}
 =\delta^\mu{}_{\nu}\pm\sqrt\lambda F^\mu{}_{\nu}
 +\lambda\,\partial^\mu\phi\,\partial_\nu\phi.
 \label{eq:Cpm-app}
\end{equation}
Equation~\eqref{eq:dbi-matrix} gives $\det C_+=\det C_-=\Delta^2$. The stress tensor in eq.~\eqref{eq:stress-dbi} may be written as
\begin{equation}
 M=\frac\Delta2(C_+^{-1}+C_-^{-1})
 =\Delta C_-^{-1}\left(\Id+\lambda\,\partial\phi\otimes\partial\phi\right)C_+^{-1}.
 \label{eq:M-factor-app}
\end{equation}
Using the matrix determinant lemma
\begin{equation}
 \det(A+uv^{\mathsf T})=\det A\left(1+v^{\mathsf T}A^{-1}u\right),
 \label{eq:detlemma-app}
\end{equation}
one immediately finds
\begin{equation}
 \det M=1+\lambda X.
\end{equation}

To obtain $E_1$, consider
\begin{equation}
 D(a)=\det\!\left(a\Id+\sqrt\lambda F
 +\lambda\,\partial\phi\otimes\partial\phi\right)
 \label{eq:Da-app}
\end{equation}
for real $a$. The four-dimensional antisymmetric matrix $F^\mu{}_{\nu}$ obeys
\begin{equation}
 (\sqrt\lambda F)^4-2\lambda S(\sqrt\lambda F)^2-\lambda^2P^2\Id=0.
 \label{eq:F-CH-app}
\end{equation}
Combining this with the matrix determinant lemma gives
\begin{equation}
 D(a)=a^4+\lambda Xa^3-2\lambda Sa^2
 +\lambda^2(Z-2SX)a-\lambda^2P^2.
 \label{eq:D-polynomial-app}
\end{equation}
The Jacobi formula,
\begin{equation}
 \frac{\dd}{\dd a}\det B(a)
 =\det B(a)\,\tr\!\left(B(a)^{-1}\frac{\dd B(a)}{\dd a}\right),
 \label{eq:jacobi-app}
\end{equation}
implies $D'(1)=\det C_+\,\tr C_+^{-1}$ and hence eq.~\eqref{eq:E1-fields}. To obtain $E_3$, we use the rank-one inverse formula
\begin{equation}
 (A+uv^{\mathsf T})^{-1}
 =A^{-1}-\frac{A^{-1}uv^{\mathsf T}A^{-1}}
 {1+v^{\mathsf T}A^{-1}u}.
 \label{eq:sherman-morrison-app}
\end{equation}
Substituting eq.~\eqref{eq:M-factor-app} into $E_3=\det M\,\tr M^{-1}$ then yields eq.~\eqref{eq:E3-fields}.

\subsection{Factorization of the characteristic polynomial}

We now derive the factorization used in sections~\ref{sec:null} and \ref{sec:constraint}. In this subsection only, write
\begin{equation}
 v_\mu\equiv\partial_\mu\phi
\end{equation}
and define
\begin{equation}
 h_{\mu\nu}=g_{\mu\nu}+\lambda v_\mu v_\nu,
 \qquad
 r\equiv\frac{\det h}{\det g}=1+\lambda X.
 \label{eq:h-r-app}
\end{equation}
We first assume $r\neq0$ and $\Delta\neq0$ so that both $h_{\mu\nu}$ and the DBI matrix are invertible. On the nonsingular real branch connected continuously to $\lambda=0$ within this parametrization, $r>0$: it starts at $r=1$ and cannot change sign without passing through the singular point $r=0$ of $h_{\mu\nu}$. The factorization obtained below can be rewritten as a polynomial identity in the original DBI variables and may then be continued algebraically to limiting loci where that identity remains defined. Such a polynomial continuation should not be confused with treating $r=0$ as part of the physical $h$-invertible branch. Introduce
\begin{equation}
 J^\mu{}_{\nu}=\sqrt\lambda\,h^{\mu\rho}F_{\rho\nu}.
 \label{eq:J-app}
\end{equation}
Then $J$ is antisymmetric with respect to $h_{\mu\nu}$ and $\G=h(\Id+J)$, so
\begin{equation}
 \Delta^2=r\det(\Id+J).
 \label{eq:Delta-J-app}
\end{equation}
The $h$-antisymmetry of $J$ implies
\begin{equation}
 (\G^{-1})^{(\mu\rho)}h_{\rho\nu}
 =\frac12\left[(\Id+J)^{-1}+(\Id-J)^{-1}\right]^\mu{}_{\nu}
 =(\Id-J^2)^{-1\,\mu}{}_{\nu}.
 \label{eq:symmetric-inverse-app}
\end{equation}
The rank-one inverse formula also gives
\begin{equation}
 (h^{-1}g)^\mu{}_{\nu}
 =\delta^\mu{}_{\nu}-\frac{\lambda}{r}v^\mu v_\nu.
 \label{eq:hinvg-app}
\end{equation}
Substituting these relations into eq.~\eqref{eq:stress-dbi} yields
\begin{equation}
 M=U\left(\Id-\frac{R}{r}\right),
 \qquad
 U\equiv\Delta(\Id-J^2)^{-1},
 \qquad
 R^\mu{}_{\nu}\equiv\lambda v^\mu v_\nu.
 \label{eq:M-U-R-app}
\end{equation}
Here $\tr R=\lambda X=r-1$.

The quadratic identity obeyed by $U$ can be established without assuming that $J$ or $U$ is diagonalizable. Since $J$ is $h$-antisymmetric in four dimensions, its characteristic polynomial is even,
\begin{equation}
 \chi_J(t)=\det(t\Id-J)=t^4+c_2t^2+c_0.
 \label{eq:J-even-char-app}
\end{equation}
Cayley--Hamilton gives
\begin{equation}
 J^4+c_2J^2+c_0\Id=0.
 \label{eq:J-CH-even-app}
\end{equation}
From $U=\Delta(\Id-J^2)^{-1}$ we have $J^2=\Id-\Delta U^{-1}$. Substitution into eq.~\eqref{eq:J-CH-even-app}, followed by multiplication by $U^2$, yields
\begin{equation}
 (1+c_2+c_0)U^2-\Delta(2+c_2)U+\Delta^2\Id=0.
 \label{eq:U-quadratic-pre-app}
\end{equation}
Furthermore,
\begin{equation}
 1+c_2+c_0=\chi_J(1)=\det(\Id-J)=\det(\Id+J)=\frac{\Delta^2}{r},
 \label{eq:J-char-at-one-app}
\end{equation}
where the last equality follows from eq.~\eqref{eq:Delta-J-app}. Therefore
\begin{equation}
 U^2-\alpha U+r\Id=0,
 \qquad
 \alpha\equiv\frac{r(2+c_2)}{\Delta}.
 \label{eq:U-quadratic-app}
\end{equation}
This derivation is insensitive to possible Jordan blocks.

The double multiplicity in the characteristic polynomial can also be shown directly at the determinant level. As a polynomial identity in $y$,
\begin{equation}
 \det(y\Id-J^2)
 =\chi_J(\sqrt y)\chi_J(-\sqrt y)
 =\left(y^2+c_2y+c_0\right)^2.
 \label{eq:J2-char-square-app}
\end{equation}
Consequently,
\begin{align}
 \det(z\Id-U)
 &=\frac{\det\!\left(z(\Id-J^2)-\Delta\Id\right)}{\det(\Id-J^2)}\nonumber\\
 &=\left[z^2-\alpha z+r\right]^2.
 \label{eq:U-char-square-app}
\end{align}
Defining $q(z)=z^2-\alpha z+r$, eq.~\eqref{eq:U-quadratic-app} then gives
\begin{equation}
 \det(z\Id-U)=q(z)^2,
 \qquad
 (z\Id-U)^{-1}U=\frac{zU-r\Id}{q(z)}.
 \label{eq:U-resolvent-app}
\end{equation}

Equation~\eqref{eq:M-U-R-app} is a rank-one update of $U$. Applying the matrix determinant lemma gives
\begin{align}
 \det(z\Id-M)
 &=q(z)^2\left[1+\frac1r\tr\!\left(R(z\Id-U)^{-1}U\right)\right]\nonumber\\
 &=q(z)\left[q(z)+\frac{\tr(RU)}{r}z-\tr R\right].
 \label{eq:char-factor-intermediate-app}
\end{align}
Using $\tr R=r-1$ and defining
\begin{equation}
 \beta\equiv\alpha-\frac{\tr(RU)}{r},
\end{equation}
we obtain
\begin{equation}
 \det(z\Id-M)
 =(z^2-\alpha z+r)(z^2-\beta z+1).
 \label{eq:char-factor-app}
\end{equation}
Since $r=E_4$ by eq.~\eqref{eq:E4}, comparison of coefficients yields
\begin{equation}
 E_1=\alpha+\beta,
 \qquad
 E_2=1+E_4+\alpha\beta,
 \qquad
 E_3=\alpha+E_4\beta.
 \label{eq:factor-coefficients-app}
\end{equation}
For $E_4\neq1$, eq.~\eqref{eq:undivided-flow} implies
\begin{equation}
 \beta=\frac{E_1-E_3}{1-E_4}=2(1+\lambda^2Q_\lambda).
 \label{eq:beta-app}
\end{equation}
At $E_4=1$, eq.~\eqref{eq:factor-coefficients-app} shows that $\alpha$ and $\beta$ are the two roots of
\begin{equation}
 t^2-E_1t+(E_2-2)=0.
\end{equation}
The field expressions in the main text give eq.~\eqref{eq:null-alpha-beta-fields}; since $Z\leq0$ for a nonzero null scalar gradient and $\Delta>0$ on the chosen DBI branch, one has $\beta\geq\alpha$. Hence
\begin{equation}
 \beta=\frac12\left[E_1+\sqrt{E_1^2-4E_2+8}\right],
\end{equation}
which yields eq.~\eqref{eq:null-E}. The discriminant is $(\alpha-\beta)^2$ on this factorized characteristic polynomial and is therefore non-negative. Eliminating $\alpha$ and $\beta$ from eq.~\eqref{eq:factor-coefficients-app} gives eq.~\eqref{eq:algebraic-relation}.

\begin{thebibliography}{99}

\bibitem{BornInfeld1934}
M.~Born and L.~Infeld,
``Foundations of the New Field Theory,''
Proc. Roy. Soc. Lond. A \textbf{144} (1934) 425--451.

\bibitem{SmirnovZamolodchikov2016}
F.~A.~Smirnov and A.~B.~Zamolodchikov,
``On space of integrable quantum field theories,''
Nucl. Phys. B \textbf{915} (2017) 363--383,
\href{https://arxiv.org/abs/1608.05499}{arXiv:1608.05499 [hep-th]}.

\bibitem{Taylor2018}
M.~Taylor,
``$T\overline T$ deformations in general dimensions,''
Adv. Theor. Math. Phys. \textbf{24} (2020) 37--63,
\href{https://arxiv.org/abs/1805.10287}{arXiv:1805.10287 [hep-th]}.

\bibitem{Bandos2020}
I.~Bandos, K.~Lechner, D.~Sorokin and P.~K.~Townsend,
``A non-linear duality-invariant conformal extension of Maxwell's equations,''
Phys. Rev. D \textbf{102} (2020) 121703,
\href{https://arxiv.org/abs/2007.09092}{arXiv:2007.09092 [hep-th]}.

\bibitem{Sorokin2021}
D.~P.~Sorokin,
``Introductory Notes on Non-linear Electrodynamics and its Applications,''
Fortsch. Phys. \textbf{70} (2022) 2200092,
\href{https://arxiv.org/abs/2112.12118}{arXiv:2112.12118 [hep-th]}.

\bibitem{Babaei2022}
H.~Babaei-Aghbolagh, K.~Babaei Velni, D.~M.~Yekta and H.~Mohammadzadeh,
``Emergence of non-linear electrodynamic theories from $T\overline T$-like deformations,''
Phys. Lett. B \textbf{829} (2022) 137079,
\href{https://arxiv.org/abs/2202.11156}{arXiv:2202.11156 [hep-th]}.

\bibitem{Ferko2022}
C.~Ferko, L.~Smith and G.~Tartaglino-Mazzucchelli,
``On Current-Squared Flows and ModMax Theories,''
SciPost Phys. \textbf{13} (2022) 012,
\href{https://arxiv.org/abs/2203.01085}{arXiv:2203.01085 [hep-th]}.

\bibitem{Conti2022}
R.~Conti, J.~Romano and R.~Tateo,
``Metric approach to a $T\overline T$-like deformation in arbitrary dimensions,''
JHEP \textbf{09} (2022) 085,
\href{https://arxiv.org/abs/2206.03415}{arXiv:2206.03415 [hep-th]}.

\bibitem{Ferko2023Birefringence}
C.~Ferko, L.~Smith and G.~Tartaglino-Mazzucchelli,
``Stress Tensor Flows, Birefringence in Non-Linear Electrodynamics, and Supersymmetry,''
SciPost Phys. \textbf{15} (2023) 198,
\href{https://arxiv.org/abs/2301.10411}{arXiv:2301.10411 [hep-th]}.

\bibitem{Ferko2023Duality}
C.~Ferko, S.~M.~Kuzenko, L.~Smith and G.~Tartaglino-Mazzucchelli,
``Duality-Invariant Non-linear Electrodynamics and Stress Tensor Flows,''
Phys. Rev. D \textbf{108} (2023) 106021,
\href{https://arxiv.org/abs/2309.04253}{arXiv:2309.04253 [hep-th]}.


\bibitem{Ferko2024DBI}
C.~Ferko, Y.~Hu, Z.~Huang, K.~Koutrolikos and G.~Tartaglino-Mazzucchelli,
``$T\overline T$-like flows and $3d$ nonlinear supersymmetry,''
SciPost Phys. \textbf{16} (2024) 038,
\href{https://arxiv.org/abs/2302.10410}{arXiv:2302.10410 [hep-th]}.

\bibitem{Blair2025}
C.~D.~A.~Blair, J.~Lahnsteiner, N.~A.~Obers and Z.~Yan,
``Matrix Theory Reloaded: A BPS Road to Holography,''
JHEP \textbf{02} (2025) 024,
\href{https://arxiv.org/abs/2410.03591}{arXiv:2410.03591 [hep-th]}.

\bibitem{Brizio2026}
N.~Brizio, M.~Kade, A.~Sfondrini and D.~P.~Sorokin,
``More on $T\overline T$-like deformations in higher dimensions,''
JHEP \textbf{06} (2026) 262,
\href{https://arxiv.org/abs/2602.24058}{arXiv:2602.24058 [hep-th]}.

\end{thebibliography}

\end{document}
